\section{Finalize without counting the padding as input}
Finalization starts with the saved used count and original byte count. Convert the
latter to a bit count with a 64-bit left shift by three. Write \texttt{80} at the next
free buffer position. If the resulting used count is greater than 56, zero the rest
of that block and compress it. Begin the next block with a used count of zero. Then
zero through position 55 and write the eight length bytes at positions 56 through 63.
Compress that final block.

The condition is stated after the marker is written. A former used count of 55 becomes
56 and still fits. A former used count of 56 becomes 57 and needs another block.
Putting the comparison on the wrong side of the marker write moves the boundary by
one byte. Test both lengths explicitly rather than inferring correctness from a short
message such as \texttt{abc}.

Serialize the full 64-bit bit count. Keeping only its low word works for small tests
but fails once the message length in bits exceeds \(2^{32}-1\). The implementation
extracts one byte at a time with doubleword logical shifts, from shift count 56 down
to zero. This also makes the byte order independent of the simulator's memory layout.
Synthetic contexts can test that serialization without allocating an enormous file;
such a test checks a supplied state and length, not the digest of a file that was
actually read.

Finally, write each state word to the output most significant byte first. For the first
word of the \texttt{abc} digest, \texttt{ba7816bf}, write bytes
\texttt{ba 78 16 bf} at increasing addresses. A direct little-endian \texttt{sw}
would write \texttt{bf 16 78 ba}. The numeric state can be correct while the
published byte sequence is wrong, so compare the output bytes as well as state words.
